\section{AGD-SDAJ verification record}\label{app:agd}

\paragraph{Gate.}
Before running our reimplementation of AGD-SDAJ~\cite{huynh2025accelerated} on the
KKT family, we required it to reproduce the authors' own published counts on
their own instances. An unverified reimplementation that is then reported as
losing would not be a defensible result.

\paragraph{Source and implementation.}
The source is Huynh and Hwang (2025), \emph{Numer.\ Linear Algebra Appl.}\
32(6), DOI \texttt{10.1002/nla.70045}. We transcribed Algorithm~1 (AGD),
Algorithm~3 (QN-SDAJ), Algorithm~4 (AGD-SDAJ) and the Section~4.3 parameters
from the full text: $c=10^{-4}$, $\rho=0.5$, $\lambda_1=\lambda_2=0.01$,
$\epsilon_l=10^{-5}$, $\epsilon_u=10^{8}$, $q=2$, $x_0=0$, the stopping rule
$\nrm{\nabla\theta}\le10^{-7}n$, and $k_{\max}=200$. The implementation is
\texttt{agd\_sdaj.jl} (standalone, used for verification) and
\texttt{agd\_sdaj\_ncm} in \texttt{bench\_sbb\_dual.jl} (instrumented, used for
the study). The instances are the Anymatrix \texttt{CORRINV} files already used
elsewhere in this project, exported with the same loading convention.

\paragraph{Equation (6), as transcribed.}
The nonmonotone line search of Algorithm~3 is a Li--Fukushima-type
rule~\cite{li2000derivative} on the residual:
\begin{equation}\label{eq:hh6}
  \nrm{\nabla\theta(x+\alpha d)}^2\ \le\ (1+\epsilon_k)\nrm{\nabla\theta(x)}^2
  -\lambda_1\nrm{\alpha\nabla\theta(x)}^2-\lambda_2\nrm{\alpha d}^2,
  \qquad \epsilon_k=\frac{1}{(k+1)^2},
\end{equation}
where, as implemented, $k$ is the index of the inner step within its outer
iteration ($\epsilon_i$ in Algorithm~\ref{alg:agd}), so that
$\epsilon_k\in\{1,\tfrac14\}$ for $q=2$. This indexing is our reading of the
text, not stated explicitly there; with it, the published counts of
\cref{tab:agd-full} are reproduced exactly on P8 and P9. The test uses
$\lambda_1=\lambda_2=0.01$, and $\alpha=\rho^{\,j}$ is backtracked from 1 over
$j=0,1,2,\dots$ (we write $j$ for the backtracking exponent, since $m$ is the
near-zero multiplicity of Definition~\ref{def:family}). The
merit function is $\tfrac12\nrm{\nabla\theta}^2$. The quasi-Newton direction
$d=-B^{-1}\nabla\theta$ need not be a descent direction for it, which is why a
nonmonotone rule is needed. The paper also confirms that $B_k$ is reset to the
identity once $x_k^{\mathrm{pre}}$ is generated, and that Step~4 discards
$x_k^{\mathrm{pre}}$ if it increases $\theta$.

Before the full text was available, we tried two guesses: a
Grippo--Lampariello--Lucidi rule on $\theta$, and a Li--Fukushima rule. The
second had the right general form but three wrong details: unsquared norms, a
missing $(1+\epsilon_k)$ slack, and $\lambda_1,\lambda_2$ attached to the
wrong terms. The guess columns of \cref{tab:agd-full} were run before the
acceleration-factor correction described below, so they are not directly
comparable with the final column. On P8 their \EVD\ counts bracketed the
published one; on P7 both were well above it, which pointed to a remaining
defect outside the line search (item~3 below).

\paragraph{Results.}
\begin{table}[h]
\centering
\caption{Outer iterations / \EVDs\ / nonmonotone backtracks, and
$\fro{G-\Xs}$, under the final transcription (\texttt{eq6}) and the two
superseded guesses (run at an earlier code state; see text).}
\label{tab:agd-full}
\small
\begin{tabular}{lllll}
\toprule
instance & published & \texttt{eq6} (final) & GLL guess & LF guess\\
\midrule
P9 \texttt{bccd16}  & 1/5/0, 29.06  & 1/5/0, 29.0563 (exact) & 1/5/0 & 1/5/0\\
P8 \texttt{cor3120} & 3/18/1, 5.44  & 3/18/1, 5.4437 (exact) & 3/17/0 & 3/19/2\\
P7 \texttt{cor1399} & 4/31/8, 21.03 & 4/32/7, 21.0339 & 6/40/3 & 5/48/17\\
\bottomrule
\end{tabular}
\end{table}
Under the final transcription, $\fro{G-\Xs}$ matched the published value to
the displayed precision on all three instances. This showed that the
objective, gradient, projection and stopping rule were correct, and that any
discrepancy lay in the amount of work per iteration.

\paragraph{Three hazards when reimplementing Algorithm~4.}
Each of the following is easy to introduce when transcribing the method from
its pseudocode. Each leaves the method converging to the correct solution
while changing the work it does, so none is visible from final errors, only
from work counts checked against the published ones. We record them for anyone
reimplementing the method. All three occurred in our own first transcription
and were found by exactly that check: the first two by accounting for the five
published \EVDs\ on P9 one at a time, the third because a gap in outer
iterations on P7 remained once the line search was correct.
\begin{enumerate}
  \item \emph{Redundant \EVD\ at an already-evaluated point.} The inner
    QN-SDAJ routine recomputed $\theta$ and $\nabla\theta$ at the point its
    caller had just evaluated. This is the same class of defect as
    \cref{sec:trap1}, and it was found in the same way. Fix: pass the
    caller's values into the inner routine.
  \item \emph{Missing early termination at the preconditioned point.} The
    note to Table~3 of~\cite{huynh2025accelerated} states that the loop also stops
    whenever $x_k^{\mathrm{pre}}$ satisfies the stopping criterion, which
    saves at least one \EVD\ (the trial at $\alpha=1$), and two when
    $\beta\ne1$. Without this check, the outer gradient step is always
    paid.
  \item \emph{Wrong acceleration factor in Algorithm~3.} The factor
    $\beta_k^{\mathrm{pre}}$ had been implemented with Algorithm~1's
    $a_k=\nrm{g}^2$. That formula is the special case for $d=-\nabla\theta$.
    In Algorithm~3, $d=-B^{-1}g$, so the general form of their Section~2.3
    (the acceleration of Andrei~\cite{andrei2006acceleration}) is needed:
    \[
      \beta=\frac{-g^{\mathsf T}d}{(g_t-g)^{\mathsf T}d}\quad\text{if the denominator is positive, else }\beta=1,
    \]
    where $g_t$ is the gradient at the accepted trial point, as in
    \cref{alg:agd}.
    This form reduces exactly to Algorithm~1's when $d=-g$, so it is
    consistent with both statements in the paper. \emph{This defect caused
    the gap in outer iterations}: correcting it reduced P7 from 6 outer
    iterations to 4.
\end{enumerate}
The cumulative effect on P7 was 46 \EVDs\ $\to$ 40 (defects 1--2) $\to$ 43
(exact equation~(6)) $\to$ 32 (correct acceleration factor), against a
published 31. (\cref{tab:agd-full} shows only the GLL-guess, LF-guess, and
final columns; the 46-\EVD\ starting point and the 43-\EVD\ intermediate,
with the exact eq.~(6) but the uncorrected acceleration factor, are not
tabled.)

\paragraph{The remaining P7 residual.}
See \cref{sec:repro}. P7 matches in outer iterations (4) and residual norm
(21.0339) and differs by one backtracking decision (7 against 8). That decision
accounts for the whole difference in \EVDs, through the early-termination
saving. We attribute it to last-bit differences between the MATLAB and Julia
eigensolvers. P8 and P9 reproduce every count exactly on the same code path.
No further parameter search was carried out.

\paragraph{Conditions of use (all met in this paper).}
(i) The code and this record are published with the results. (ii) The P7
residual is disclosed. (iii) Published counts are printed beside local counts
wherever the comparison appears (\cref{tab:agd-verify,tab:agd-full}).

\paragraph{Checks against other published results.}
Beyond the AGD-SDAJ gate, every solver was run on invalid correlation matrices
with published solutions or work counts
(\cref{tab:external}). The matrices are those of the \texttt{CORRINV}
collection~\cite{higham2022anymatrix}, taken from its repository rather than
retyped. On the ten small matrices (\texttt{high02}, \texttt{tec03},
\texttt{bhwi01}, \texttt{mmb13}, \texttt{fing97}, \texttt{beyu11},
\texttt{tyda99r1}--\texttt{r3}, \texttt{usgs13}), run at tolerance $10^{-10}$,
the five solvers of the original comparison return valid correlation
matrices (Anderson-APM, added later, is checked in \cref{app:anderson}) whose
distances
$\fro{G-X}$ agree to at least eight decimal places and whose matrices agree
with Newton-SIN-BH's to within $2\times10^{-10}$, except unmodified AGD-SDAJ on
\texttt{mmb13} ($8.7\times10^{-10}$), which stops at its budget. The three
large matrices (\texttt{cor1399}, \texttt{cor3120} and \texttt{bccd16}) were
compared only at the published tolerances.

\begin{table}[h]
\centering
\caption{Work counts reproduced from the literature. Dykstra-APM uses the
stopping rule of the released code \texttt{nearcorr\_new.m}
of~\cite{higham2016anderson}, tolerance $n\,\mathrm{eps}$; run for the same
number of iterations, our implementation returns the same iterate, bit for
bit on nine matrices and to $1.4\times10^{-15}$ on \texttt{mmb13}. Newton-SIN-BH is compared with Tables 6.5 and 6.6
of~\cite{borsdorf2007newton}.}
\label{tab:external}
\small
\begin{tabular}{llll}
\toprule
solver & source & matrix & published / ours\\
\midrule
Dykstra-APM & \cite{higham2016anderson}, Table 1 & \texttt{tec03}, $n=4$ & 39 / 39 iterations\\
            & & \texttt{bhwi01}, $n=5$ & 27 / 27\\
            & & \texttt{mmb13}, $n=6$ & 801 / 801\\
            & & \texttt{fing97}, $n=7$ & 33 / 34$^{a}$\\
            & \cite{higham2016anderson}, Table 4 & \texttt{usgs13}, $n=94$ & 18 / 18\\
\midrule
Newton-SIN-BH & \cite{borsdorf2007newton} & \texttt{cor1399}, $10^{-7}n$ & 5 / 5 iterations\\
              & & \texttt{cor1399}, $n\,\mathrm{eps}$ & 7 / 7; $\fro{G-X}=21.033911$ / 21.033911\\
              & & \texttt{cor3120}, $10^{-7}n$ & 4 / 4\\
              & & \texttt{cor3120}, $n\,\mathrm{eps}$ & 6 / not reached$^{b}$; 5.443751 / 5.443751\\
\bottomrule
\end{tabular}
\end{table}
{\small $^{a}$At iteration 33 the stopping quantity is $1.12$ times the
tolerance $7\,\mathrm{eps}$, so the extra iteration is decided in the last
bits. $^{b}$The residual floor near $10^{-12}$ lies above
$n\,\mathrm{eps}=6.9\times10^{-13}$ with our eigensolver; the solution itself
matches to all published digits.\par}

On three of these matrices unmodified AGD-SDAJ shows the effect of
\cref{sec:trap2}: on \texttt{mmb13} it stalls (budget of 200\,000 \EVDs\
exhausted, against 210 for AGD-SDAJ-BH), and on \texttt{tyda99r2} and
\texttt{bhwi01} it converges, but slowly (549 against 35, and 66 against 21).

\paragraph{The Borsdorf--Higham adaptation.}
AGD-SDAJ-BH changes only the monotone Armijo test of Algorithm~1, line~5
(\cref{alg:agd}). When
$|\theta_t-\theta|<\gamma u(1+|\theta_t|+|\theta|)$ with $\gamma=100$ and
$u=2^{-53}$, a trial that reduces $\nrm{\nabla\theta}$ is accepted (weaker than
the Newton safeguard's $(1-\mubh)$ reduction by design: near the precision floor
a gradient step only inches the gradient down, so the stronger test never
fires). The
nonmonotone rule~\eqref{eq:hh6} is unchanged. This adaptation is ours.
